% !PRINTED 445
\weberpage{445}
be divisible by $p$, because both $\nu - \nu'$ and $\nu + \nu'$ are smaller than
$p$. Among the numbers

\weq{s}{1, 2, 3 \ldots p - 1}

half therefore occur among the residues of $x^2$, and half do not.
The former numbers, and all those congruent to them with respect to $p$,
are called the \gesperrt{quadratic residues of $p$}, the others the
\gesperrt{quadratic non-residues of $p$}.

Now if $m$ is a quadratic residue, the congruence
\[
x^2 \equiv m \ (\mathrm{mod}\ p)
\]
is possible, and from 7. and 8. it follows that
\[
\left(\frac{m}{p}\right) = \left(\frac{x^2}{p}\right) = \left(\frac{x}{p}\right)^2 = +1.
\]
Thus $\left(\dfrac{m}{p}\right) = +1$ when $m$ is a quadratic residue of
$p$.

But in \S.\,135, 1. the theorem was stated that, if for $m$ the
sequence of numbers 1, 2, \ldots p - 1 is put,
$\left(\dfrac{m}{p}\right)$ has the positive sign just as often as the
negative one; and from this one sees that the negative sign holds
when $m$ is a quadratic non-residue. It therefore follows:

11. \gesperrt{If $p$ is an odd prime and $m$ is not divisible
by $p$, then $\left(\dfrac{m}{p}\right) = +1$ or $= -1$,
according as $m$ is a quadratic residue or a quadratic non-residue of
$p$.}
In this all the principal theorems on the quadratic residues
are contained, e.\,g. the theorem:

The product of two quadratic residues or of two non-residues
is a residue; the product of a residue and a non-residue is
a non-residue.

\begin{center}
\rule{0.25\linewidth}{0.4pt}
\end{center}
